% GENERATED FILE. Edit canonical lessons, then run npm run generate:books.
% canonical-source-hash: fnv1a64:e1f7258f7df3b3b1
% canonical-lessons: ZH-C37-zhe, ZH-W37-zhe, ZH-C37-zhege, ZH-C37-na, ZH-W37-na, ZH-C37-nage-which, ZH-C37-nage

\chapter{This One and That One}
\label{ch:zh-this-one-and-that-one}
\hlchaptermodality{37}

\begin{chapteropening}
\textbf{By the end of this chapter:} \emph{I can point things out as this one, that one and which one.}

Complete the last lesson of chapter 37: I can point things out as this one, that one and which one.
\end{chapteropening}

\section[\texorpdfstring{\zh{这} (zhè)}{这 (zhè)}]{\zh{这} (zhè) — this}

\label{lesson:ZH-C37-zhe}



\begin{quote}
Before the new one: say the Mandarin for Beijing, then the Mandarin for which.
\end{quote}

\subsection*{You'll want to know: \zh{这}}
\textbf{\zh{这}} — \emph{zhè} — \textquotedblleft{}this\textquotedblright{}.

\begin{quote}
\textbf{\zh{这是我的书}} \emph{zhè shì wǒ de shū} — This is my book.
\end{quote}

The new character, in pieces: \zh{这} — \zh{文}, from \zh{中文}, and a walking piece that runs under it from the left. The next lesson puts it in your hand.

\begin{cousinweb}
\textbf{\zh{这}} (\emph{zhè}) is written with an old character that the Guangyun records as \emph{ngjenH} (Guangyun, in Baxter's transcription), glossed 'to meet'. That is a different word: \emph{ngjenH} would not become zhè. How the character came to write the word for this is not settled.
\end{cousinweb}

\begin{grammarlens}[title={What you've built}]
One more word for everyday answers.
\end{grammarlens}

\subsection*{Your turn}
\begin{itemize}
\raggedright
  \item \emph{Say it:} \emph{zhè}
  \item \emph{Say it:} \emph{zhè}, once more
  \item \emph{Say it:} \emph{zhè shì wǒ de shū} — This is my book.
  \item \emph{Recall:} say the Mandarin for Beijing, then the Mandarin for which, then say \emph{zhè} again
\end{itemize}

\begin{tcolorbox}[breakable,colback=teal!4,colframe=teal!35!black,title={Before you move on}]
What does \zh{这} mean? (\textquotedblleft{}This\textquotedblright{}.) Say it once more.
\end{tcolorbox}
\section[\texorpdfstring{\zh{这} (zhè)}{这 (zhè)}]{\zh{这} — \zh{文}, then the walking piece}

\label{lesson:ZH-W37-zhe}



\begin{quote}
The word first, then the character inside it.

\begin{itemize}
\raggedright
  \item \emph{Recall:} say \emph{zhè} — \textbf{R1}, one lesson back, before it decays
\end{itemize}

\begin{quote}
\zh{文} \zh{书}
\end{quote}

\emph{Read it:} each one aloud, with nothing said first — \emph{wén}, \emph{shū}
\end{quote}

\subsection*{Script — \zh{文}, then the walking piece}
\hlblockfigure{\detokenize{figures/ZH-W37-zhe-filmstrip.pdf}}{How \zh{这} is written, stroke by stroke}

\begin{quote}
\zh{这}
\end{quote}

\textbf{Two pieces, and the one underneath goes last.} The walking piece wraps the bottom-left, but it is written after what it carries.

\textbf{The top-right piece, \zh{文}, four strokes}, as in \zh{中文}: a dot, a horizontal, a left-falling stroke and a right-falling stroke crossing it.

\textbf{The walking piece, three strokes.} A dot. A stroke that runs across, folds down and folds again, ending in a short fall. Then a long flat sweep along the bottom, under \zh{文}, out to the right.

4 + 3 = seven.

\subsection*{Your turn}
\begin{itemize}
\raggedright
  \item \emph{Write it:} \zh{这}, all seven, with the model in front of you
  \item \emph{Cover:} write it again from memory, and count the pieces, not the strokes
\end{itemize}

\begin{tcolorbox}[breakable,colback=teal!4,colframe=teal!35!black,title={Before you move on}]
How many pieces in \zh{这}? (\textbf{Two} — \zh{文}, then the walking piece.) Which goes first? (\textbf{\zh{文}}, even though the walking piece is on the left.) Which stroke is last? (\textbf{The long sweep along the bottom}.)

Source: \href{https://raw.githubusercontent.com/chanind/hanzi-writer-data/68d10a4b21150cae5e1ebbd223eed289cf32d90c/data/\%E8\%BF\%99.json}{Hanzi Writer Data \zh{这}, pinned stroke paths and medians}.
\end{tcolorbox}
\section[\texorpdfstring{\zh{这个} (zhège)}{这个 (zhège)}]{\zh{这个} (zhège) — this one}

\label{lesson:ZH-C37-zhege}



\begin{quote}
Before the new one: say the Mandarin for which, then the Mandarin for this.
\end{quote}

\subsection*{You'll want to know: \zh{这个}}
\textbf{\zh{这个}} — \emph{zhège} — \textquotedblleft{}this one\textquotedblright{}.

\begin{quote}
\textbf{\zh{这个人叫什么名字}} \emph{zhège rén jiào shénme míngzi} — What's this person's name?
\end{quote}

\begin{cousinweb}
\textbf{\zh{这个}} (\emph{zhège}) is built from \zh{这} 'this' and the measure word \zh{个}.
\end{cousinweb}

\begin{grammarlens}[title={What you've built}]
One more word for everyday answers.
\end{grammarlens}

\subsection*{Your turn}
\begin{itemize}
\raggedright
  \item \emph{Say it:} \emph{zhège}
  \item \emph{Say it:} \emph{zhège}, once more
  \item \emph{Say it:} \emph{zhège rén jiào shénme míngzi} — What's this person's name?
  \item \emph{Recall:} say the Mandarin for which, then the Mandarin for this, then say \emph{zhège} again
\end{itemize}

\begin{tcolorbox}[breakable,colback=teal!4,colframe=teal!35!black,title={Before you move on}]
What does \zh{这个} mean? (\textquotedblleft{}This one\textquotedblright{}.) Say it once more.
\end{tcolorbox}
\section[\texorpdfstring{\zh{那} (nà)}{那 (nà)}]{\zh{那} (nà) — that}

\label{lesson:ZH-C37-na}



\begin{quote}
Before the new one: say the Mandarin for this, then the Mandarin for this one.
\end{quote}

\subsection*{You'll want to know: \zh{那}}
\textbf{\zh{那}} — \emph{nà} — \textquotedblleft{}that\textquotedblright{}.

\begin{quote}
\textbf{\zh{那是我的家}} \emph{nà shì wǒ de jiā} — That's my home.
\end{quote}

The new character, in pieces: \zh{那} — a four-stroke piece on the left, and on the right a two-stroke piece shaped like an ear. The next lesson puts it in your hand.

\begin{cousinweb}
\textbf{\zh{那}} (\emph{nà}) is an old character: the Guangyun records it as \emph{na} (Guangyun, in Baxter's transcription), glossed 'how, what'. The sense 'that' is not given there, so how it came to mean that is not settled.
\end{cousinweb}

\begin{grammarlens}[title={What you've built}]
One more word for everyday answers.
\end{grammarlens}

\subsection*{Your turn}
\begin{itemize}
\raggedright
  \item \emph{Say it:} \emph{nà}
  \item \emph{Say it:} \emph{nà}, once more
  \item \emph{Say it:} \emph{nà shì wǒ de jiā} — That's my home.
  \item \emph{Recall:} say the Mandarin for this, then the Mandarin for this one, then say \emph{nà} again
\end{itemize}

\begin{tcolorbox}[breakable,colback=teal!4,colframe=teal!35!black,title={Before you move on}]
What does \zh{那} mean? (\textquotedblleft{}That\textquotedblright{}.) Say it once more.
\end{tcolorbox}
\section[\texorpdfstring{\zh{那} (nà)}{那 (nà)}]{\zh{那} — four strokes, then an ear}

\label{lesson:ZH-W37-na}



\begin{quote}
The word first, then the character inside it.

\begin{itemize}
\raggedright
  \item \emph{Recall:} say \emph{nà} — \textbf{R1}, one lesson back, before it decays
\end{itemize}

\begin{quote}
\zh{这} \zh{二}
\end{quote}

\emph{Read it:} each one aloud, with nothing said first — \emph{zhè}, \emph{èr}
\end{quote}

\subsection*{Script — four strokes, then an ear}
\hlblockfigure{\detokenize{figures/ZH-W37-na-filmstrip.pdf}}{How \zh{那} is written, stroke by stroke}

\begin{quote}
\zh{那}
\end{quote}

\textbf{Two pieces, left before right.}

\textbf{The left piece, four strokes.} A horizontal that turns down and hooks. Two short horizontals inside it. Then a left-falling stroke from the top, down through them all.

\textbf{The ear, two strokes.} One stroke that runs right, loops down twice like a folded ear and hooks. Then a long vertical down its left edge.

4 + 2 = six.

\subsection*{Your turn}
\begin{itemize}
\raggedright
  \item \emph{Write it:} \zh{那}, all six, with the model in front of you
  \item \emph{Cover:} write it again from memory, and count the pieces, not the strokes
\end{itemize}

\begin{tcolorbox}[breakable,colback=teal!4,colframe=teal!35!black,title={Before you move on}]
How many pieces in \zh{那}? (\textbf{Two} — four strokes, then the ear.) Which stroke of the left piece is last? (\textbf{The left-falling stroke through it all}.) Which is the last stroke of the character? (\textbf{The long vertical}.)

Source: \href{https://raw.githubusercontent.com/chanind/hanzi-writer-data/68d10a4b21150cae5e1ebbd223eed289cf32d90c/data/\%E9\%82\%A3.json}{Hanzi Writer Data \zh{那}, pinned stroke paths and medians}.
\end{tcolorbox}
\section[nǎge]{nǎge — which one}

\label{lesson:ZH-C37-nage-which}



\begin{quote}
Before the new one: say the Mandarin for this one, then the Mandarin for that.
\end{quote}

\subsection*{You'll want to know: nǎge}
\emph{nǎge} — \textquotedblleft{}which one\textquotedblright{}.

\begin{quote}
\emph{nǐ xǐhuan nǎge} — Which one do you like?
\end{quote}

Its characters come later; for now it is a word for the ear.

\begin{cousinweb}
\emph{nǎge} is built from nǎ 'which' and the measure word \zh{个}.
\end{cousinweb}

\begin{grammarlens}[title={What you've built}]
One more word for everyday answers.
\end{grammarlens}

\subsection*{Your turn}
\begin{itemize}
\raggedright
  \item \emph{Say it:} \emph{nǎge}
  \item \emph{Say it:} \emph{nǎge}, once more
  \item \emph{Say it:} \emph{nǐ xǐhuan nǎge} — Which one do you like?
  \item \emph{Recall:} say the Mandarin for this one, then the Mandarin for that, then say \emph{nǎge} again
\end{itemize}

\begin{tcolorbox}[breakable,colback=teal!4,colframe=teal!35!black,title={Before you move on}]
What does nǎge mean? (\textquotedblleft{}Which one\textquotedblright{}.) Say it once more.
\end{tcolorbox}
\section[\texorpdfstring{\zh{那个} (nàge)}{那个 (nàge)}]{\zh{那个} (nàge) — that one}

\label{lesson:ZH-C37-nage}



\begin{quote}
Before the new one: say the Mandarin for that, then the Mandarin for which one.
\end{quote}

\subsection*{You'll want to know: \zh{那个}}
\textbf{\zh{那个}} — \emph{nàge} — \textquotedblleft{}that one\textquotedblright{}.

\begin{quote}
\textbf{\zh{我喜欢那个}} \emph{wǒ xǐhuan nàge} — I like that one.
\end{quote}

\begin{cousinweb}
\textbf{\zh{那个}} (\emph{nàge}) is built from \zh{那} 'that' and the measure word \zh{个}.
\end{cousinweb}

\begin{grammarlens}[title={What you've built}]
One more word for everyday answers.
\end{grammarlens}

\subsection*{Your turn}
\begin{itemize}
\raggedright
  \item \emph{Say it:} \emph{nàge}
  \item \emph{Say it:} \emph{nàge}, once more
  \item \emph{Say it:} \emph{wǒ xǐhuan nàge} — I like that one.
  \item \emph{Recall:} say the Mandarin for that, then the Mandarin for which one, then say \emph{nàge} again
\end{itemize}

\begin{tcolorbox}[breakable,colback=teal!4,colframe=teal!35!black,title={Before you move on}]
What does \zh{那个} mean? (\textquotedblleft{}That one\textquotedblright{}.) Say it once more.
\end{tcolorbox}
